\section{Conclusion}
In this work, we introduced Think-on-Graph 3.0, a novel framework that fundamentally rethinks the paradigm of RAG for complex reasoning. 
By proposing the Multi-Agent Context Evolution Retrieval (MACER) mechanism and a dynamic Chunk-Triplets-Community heterogeneous graph architecture, we address critical limitations in both existing graph-based RAG methods and knowledge-graph-dependent approaches.
Our comprehensive experimental evaluation demonstrates that ToG-3 achieves state-of-the-art performance across multiple challenging benchmarks. 
% The framework's core innovation is its dual-evolving mechanism—comprising Evolving Query and Evolving Subgraph—which dynamically refines both the query representation and the underlying graph structure throughout the reasoning process. This iterative co-evolution enables deep, multi-hop inference while preserving broad coverage across complex queries.
This adaptive capability proves particularly valuable for overcoming the quality constraints of static graph construction and the domain limitations of pre-existing knowledge bases.
The framework's ability to work with light LLMs also opens possibilities for more efficient and deployable AI systems.
% We believe the principles established in ToG-3—dynamic graph evolution, multi-agent collaboration, and heterogeneous knowledge integration—provide a foundation for the next generation of RAG systems. 
% that can truly bridge the gap between language understanding and knowledge reasoning.